\chapter{Nội dung thực tập tại phòng thí nghiệm}
\section{Vị trí thực tập}
\subsection{Tên vị trí}
\textbf{Thực tập sinh Kiểm định chất lượng sản phẩm -- phòng thí nghiệm}

\subsection{Mô tả công việc}
Mục tiêu công việc: Hỗ trợ phòng thí nghiệm trong hoạt động kiểm tra, phân tích và đánh giá chất lượng nguyên liệu, bán thành phẩm và thành phẩm phân bón; đồng thời hỗ trợ ghi nhận, quản lý kết quả kiểm nghiệm phục vụ công tác kiểm soát chất lượng của công ty \cite{company-internal-documents}.

Nhiệm vụ chính:
\begin{itemize}
  \item Tiếp nhận, phân loại, mã hóa và chuẩn bị mẫu nguyên liệu, bán thành phẩm và thành phẩm trước khi phân tích.
  \item Thực hiện các công việc xử lý mẫu như nghiền, cân, hòa tan, pha loãng và chuẩn bị dung dịch theo hướng dẫn.
  \item Chuẩn bị hóa chất, thuốc thử, dụng cụ và thiết bị cần thiết cho từng phép kiểm nghiệm.
  \item Hỗ trợ thực hiện các phép phân tích và kiểm tra chỉ tiêu chất lượng của sản phẩm phân bón theo quy trình của phòng thí nghiệm.
  \item Ghi chép, tổng hợp và lưu trữ kết quả phân tích; hỗ trợ đối chiếu kết quả với chỉ tiêu kỹ thuật của từng sản phẩm.
  \item Theo dõi và lưu trữ mẫu phục vụ việc kiểm tra, đối chứng hoặc truy xuất khi cần thiết.
  \item Vệ sinh dụng cụ, thiết bị và khu vực làm việc sau khi hoàn thành quá trình phân tích.
  \item Sắp xếp hóa chất, dụng cụ và mẫu đúng vị trí quy định, bảo đảm điều kiện bảo quản phù hợp.
  \item Tuân thủ các quy định về an toàn hóa chất, an toàn thiết bị và nội quy phòng thí nghiệm.
  \item Hỗ trợ nhân viên phòng thí nghiệm trong các công việc kiểm soát chất lượng và nghiên cứu thử nghiệm khác khi được phân công.
\end{itemize}

\section{Quy trình lấy mẫu kiểm nghiệm}
Quá trình lấy mẫu được thực hiện tại ba vị trí chính trên dây chuyền, bao gồm nguyên liệu đầu vào, bán thành phẩm trong quá trình sản xuất và thành phẩm cuối dây chuyền\cite{company-internal-documents}. 
Trong cả ba giai đoạn, quá trình lấy mẫu cần tuân thủ một số nguyên tắc cơ bản \cite{company-internal-documents}:
\begin{itemize}
  \item Mẫu phải có tính đại diện cho lô nguyên liệu hoặc sản phẩm cần kiểm tra.
  \item Dụng cụ lấy và chứa mẫu phải sạch, khô, không làm thay đổi thành phần của mẫu.
  \item Không để lẫn mẫu giữa các lô hoặc các sản phẩm khác nhau.
  \item Sau mỗi lần lấy mẫu cần vệ sinh dụng cụ trước khi sử dụng cho mẫu tiếp theo.
  \item Mẫu phải được đóng kín và ghi nhãn ngay sau khi lấy.
  \item Nhãn mẫu cần thể hiện tối thiểu tên mẫu, mã lô, vị trí lấy mẫu, ngày giờ lấy mẫu và người thực hiện.
  \item Mẫu phải được chuyển về phòng thí nghiệm trong điều kiện phù hợp và hạn chế tối đa sự thay đổi trong quá trình vận chuyển.
  \item Số lượng điểm lấy mẫu, lượng mẫu và phương pháp chia mẫu được thực hiện theo quy trình kiểm soát chất lượng áp dụng cho từng loại nguyên liệu và sản phẩm của công ty.
  \item Các mẫu sau khi lấy được bảo quản trong điều kiện phù hợp để tránh biến đổi thành phần trước khi phân tích.
\end{itemize}

\subsection{Lấy mẫu nguyên liệu đầu vào}
Mẫu nguyên liệu được lấy tại đầu dây chuyền sản xuất, trước khi nguyên liệu được đưa vào phối trộn. 
Việc lấy mẫu nhằm kiểm tra chất lượng từng lô nguyên liệu và phát hiện các trường hợp không phù hợp trước khi đưa vào sản xuất. 

Đối với mỗi lô nguyên liệu nhập về, tiến hành lấy mẫu ít nhất một lần cho mỗi lô hoặc mỗi chuyến nguyên liệu. 
Mẫu được lấy tại nhiều vị trí khác nhau trong lô, với mỗi mẫu đơn khoảng 100--200 g \cite{company-internal-documents}.
Các mẫu đơn sau khi lấy được gộp lại, trộn đồng nhất và chuyển về phòng thí nghiệm. 
Lập mẫu đại diện hoàn chỉnh cho lô, mẫu chung được rút gọn đến trên 1,5 kg, sau đó chia thành ba phần, mỗi phần tối thiểu 0,5 kg để phục vụ phân tích, lưu mẫu và đối chứng \cite{tcvn9486-sampling}.

Mẫu nguyên liệu sau khi lấy phải được ghi đầy đủ thông tin như tên sản phẩm, công thức, mã lô, ngày sản xuất, thời gian lấy mẫu và người thực hiện. 
Mẫu lưu được đóng kín và bảo quản tại khu vực lưu mẫu của phòng thí nghiệm để phục vụ việc truy xuất hoặc kiểm tra lại khi cần thiết \cite{company-internal-documents}.

\begin{figure}[H]
  \centering
  \includegraphics[width=0.72\textwidth,height=0.56\textheight,keepaspectratio]{lay-mau-nguyen-lieu.jpg}
  \caption{Mẫu nguyên liệu đầu vào sau khi lấy và ghi nhãn.}
  \label{fig:lay-mau-nguyen-lieu}
\end{figure}


\subsection{Lấy mẫu bán thành phẩm}
Mẫu bán thành phẩm được lấy tại giữa dây chuyền sản xuất, sau khi các nguyên liệu đã trải qua công đoạn phối trộn nhưng trước khi sản phẩm hoàn thiện. 
Mục đích của bước này là kiểm tra mức độ đồng đều của quá trình phối trộn và phát hiện sớm sự sai lệch về thành phần \cite{company-internal-documents}.

Sau khi dây chuyền bắt đầu hoạt động và đạt trạng thái tương đối ổn định, tiến hành lấy mẫu đầu tiên. 
Trong quá trình sản xuất liên tục, mẫu bán thành phẩm được lấy với tần suất khoảng 30 phút/lần, mỗi lần khoảng 100--200g sao cho đảm bảo tính đại diện cho lô sản xuất và mẫu chung đảm bảo trên 1,5 kg sau khi gộp lại để phục vụ mục đích phân tích, lưu mẫu và đối chứng \cite{company-internal-documents}.

Khi thay đổi công thức sản phẩm, thay đổi nguồn nguyên liệu hoặc dây chuyền dừng và khởi động lại, chu kỳ lấy mẫu được tính lại từ đầu. 

Mẫu bán thành phẩm sau khi lấy phải được ghi đầy đủ thông tin như tên sản phẩm, công thức, mã lô, ngày sản xuất, thời gian lấy mẫu và người thực hiện. 
Mẫu lưu được đóng kín và bảo quản tại khu vực lưu mẫu của phòng thí nghiệm để phục vụ việc truy xuất hoặc kiểm tra lại khi cần thiết \cite{company-internal-documents}.

\subsection{Lấy mẫu thành phẩm}
Mẫu thành phẩm được lấy tại cuối dây chuyền, sau khi quá trình phối trộn và xử lý hoàn tất, trước hoặc trong quá trình đóng bao. 
Đây là mẫu quan trọng để đánh giá chất lượng cuối cùng của lô sản xuất. 

Trong quá trình thành phẩm ra khỏi dây chuyền, SOP nội bộ đề xuất lấy mẫu với tần suất khoảng 30 phút/lần ở các vị trí khác nhau như ngay khi hoàn tất quy trình phối trộn hoặc trong quá trình đóng bao. 
Mỗi lần thường lấy khoảng 100--200 g, và duy trì xuyên suốt thời gian sản xuất của lô sao cho đảm bảo tính đại diện cho lô sản xuất và mẫu chung đảm bảo trên 1,5 kg sau khi gộp lại để phục vụ mục đích phân tích, lưu mẫu và đối chứng \cite{company-internal-documents}.
Đối với lô kéo dài nhiều giờ, tiếp tục lấy mẫu định kỳ 30 phút/lần cho đến khi kết thúc lô. Sau khi hoàn thành lô sản xuất, toàn bộ mẫu đơn được gộp và trộn đều để tạo thành mẫu chung. Theo TCVN 9486:2018, mẫu chung dạng rắn được rút gọn bằng phương pháp chia mẫu đến khi còn trên 1,5 kg, sau đó chia thành ba đơn vị mẫu, mỗi đơn vị tối thiểu 0,5 kg \cite{tcvn9486-sampling}.

Mẫu thành phẩm sau khi lấy phải được ghi đầy đủ thông tin như tên sản phẩm, công thức, mã lô, ngày sản xuất, thời gian lấy mẫu và người thực hiện. 
Mẫu lưu được đóng kín và bảo quản tại khu vực lưu mẫu của phòng thí nghiệm để phục vụ việc truy xuất hoặc kiểm tra lại khi cần thiết \cite{company-internal-documents}.

\begin{figure}[H]
  \centering
  \includegraphics[width=0.72\textwidth,height=0.56\textheight,keepaspectratio]{lay-mau-thanh-pham.jpg}
  \caption{Các mẫu thành phẩm sau khi lấy, đóng gói và ghi phiếu lấy mẫu.}
  \label{fig:lay-mau-thanh-pham}
\end{figure}


\begin{table}[H]
  \centering
  \caption{Tóm tắt tần suất lấy mẫu áp dụng tại phòng thí nghiệm.}
  \label{tab:sampling-frequency}
  \begin{tabular}{p{0.23\textwidth} p{0.14\textwidth} p{0.28\textwidth} p{0.22\textwidth}}
    \toprule
    Mẫu kiểm nghiệm & Lượng lấy mỗi lần & Tần suất đề xuất & Mục đích chính \\
    \midrule
    Nguyên liệu đầu vào & 100--200 g/lần & Mỗi lô/chuyến: tại nhiều vị trí khác nhau & Kiểm tra chất lượng nguyên liệu đầu vào\\
    Bán thành phẩm & 100--200 g/lần & Khoảng 30 phút/lần & Theo dõi độ đồng đều, phát hiện sai lệch sớm \\
    Thành phẩm & 100--200 g/lần & Khoảng 30 phút/lần; tại các vị trí khác nhau & Đánh giá chất lượng toàn lô \\
    Mẫu đại diện cuối cùng & Trên 1,5 kg & Sau khi kết thúc lô & Chia ba phần, mỗi phần tối thiểu 0,5 kg \\
    \bottomrule
  \end{tabular}
  \par\smallskip
  \small Nguồn: SOP nội bộ của công ty và TCVN 9486:2018 \cite{company-internal-document,tcvn9486-sampling}.
\end{table}


\section{Quy trình kiểm tra chất lượng}
Quy trình kiểm tra chất lượng tại phòng thí nghiệm được thực hiện đối với nguyên liệu đầu vào, bán thành phẩm và thành phẩm nhằm đánh giá mức độ phù hợp của lô hàng với công thức sản phẩm và chỉ tiêu kỹ thuật áp dụng. 
Trong phạm vi thực tập, các chỉ tiêu cơ bản được theo dõi gồm đạm tổng số, kali hữu hiệu quy đổi theo \(\mathrm{K_2O}\), lân hữu hiệu quy đổi theo \(\mathrm{P_2O_5}\), carbon tổng và độ ẩm. Việc lựa chọn phương pháp phụ thuộc vào nền mẫu, dải nồng độ, thiết bị hiện có và quy trình nội bộ của công ty \cite{company-internal-documents,qcvn-fertilizer-quality}.

\subsection{Nguyên tắc chung}
\begin{itemize}
  \item Trước khi phân tích, mẫu được kiểm tra nhãn, mã lô và tình trạng bao gói; sau đó được nghiền, trộn đồng nhất và rút gọn theo quy định để bảo đảm tính đại diện. 
  \item Dụng cụ, hóa chất, thiết bị đo và dung dịch chuẩn phải phù hợp với phép thử. 
  \item Các thiết bị như cân phân tích, tủ sấy hoặc máy đo ngọn lửa được kiểm tra tình trạng hoạt động và hiệu chuẩn theo kế hoạch của phòng thí nghiệm trước khi sử dụng \cite{company-internal-documents}.
  \item Mỗi phép thử cần có mẫu trắng, mẫu chuẩn hoặc mẫu kiểm soát chất lượng khi phù hợp. 
  \item Kết quả được ghi trực tiếp vào biểu mẫu, kiểm tra tính hợp lý và đối chiếu với giới hạn kỹ thuật của sản phẩm. 
  \item Khi có dấu hiệu bất thường, nhân viên phụ trách xem xét việc phân tích lặp lại, pha loãng lại mẫu hoặc kiểm tra chéo bằng mẫu lưu trước khi đưa ra kết luận \cite{company-internal-documents}.
\end{itemize}

\subsection{Xác định đạm tổng số}
\subsubsection{Nguyên tắc}
Tại phòng thí nghiệm của Công ty TNHH Quốc tế Khánh Sinh, phương pháp Kjeldahl kết hợp chưng cất được thực hiện để xác định hàm lượng đạm tổng số trong mẫu phân bón.
Cụ thể, các hợp chất chứa nitơ trong mẫu được vô cơ hóa bằng sulfuric acid (\(\mathrm{H_2SO_4}\)) đậm đặc trong môi trường xúc tác để chuyển hóa thành ammonium sulfate.
\[
  \text{Hợp chất chứa N}
  \xrightarrow[\Delta]{\mathrm{H_2SO_4}}
  \mathrm{(NH_4)_2SO_4 + CO_2 + H_2O + \text{sản phẩm khác}}
\]

Sau đó, dung dịch được kiềm hóa bằng \(\mathrm{NaOH}\) để giải phóng khí ammonia (\(\mathrm{NH_3}\)).
\[
  \mathrm{NH_4^+ + OH^- \longrightarrow NH_3 \uparrow + H_2O}
\]
Khí ammonia được dẫn vào bình hứng chứa dung dịch boric acid (\(\mathrm{H_3BO_3}\)) để hấp thụ và tạo thành hệ ammonium tetrahydroxyborate.
\[
  \mathrm{NH_3 + H_3BO_3 + H_2O \longrightarrow NH_4^+ + B(OH)_4^-}
\]

Cuối cùng, lượng nitơ trong mẫu được xác định bằng cách chuẩn độ dung dịch ammonium tetrahydroxyborate với dung dịch sulfuric acid chuẩn (\(\mathrm{H_2SO_4}\)) 0,1 N, sử dụng chỉ thị màu hỗn hợp methyl red - bromocresol green để xác định điểm kết thúc của phản ứng.
\[
  \mathrm{2NH_4^+ + 2B(OH)_4^- + H_2SO_4 \longrightarrow (NH_4)_2SO_4 + 2H_3BO_3 + 2H_2O}
\]
Phản ứng trung hoà được biểu thị bằng sự thay đổi màu sắc từ xanh thành đỏ của dung dịch.

\subsubsection{Thiết bị, hóa chất và thuốc thử}
+ Hóa chất và thuốc thử.
\begin{itemize}
  \item Sulfuric acid (\(\mathrm{H_2SO_4}\)) đậm đặc (d = 1,84).
  \item Sulfuric acid chuẩn (\(\mathrm{H_2SO_4}\)) 0,1 N
  \item Sodium hydroxide (\(\mathrm{NaOH}\)) tinh thể
  \item Methyl red (\(\mathrm{C_{15}H_{15}N_3O_2}\)) tinh thể
  \item Bromocresol green (\(\mathrm{C_{21}H_{14}Br_4O_5S}\)) tinh thể
  \item Potassium sulfate (\(\mathrm{K_2SO_4}\)) tinh thể
  \item Copper(II) sulfate dihydrate (\(\mathrm{CuSO_4\cdot 2H_2O}\)) tinh thể
  \item Boric acid (\(\mathrm{H_3BO_3}\)) tinh thể
  \item Giấy chỉ thị pH
  \item Ammonium sulfate (\(\mathrm{(NH_4)_2SO_4}\))
  \item Phenolphthalein (\(\mathrm{C_{20}H_{12}O_4}\)) tinh thể
\end{itemize}

+ Thiết bị và dụng cụ
\begin{itemize}
  \item Các dụng cụ thông thường trong phòng thí nghiệm.
  \item Cân phân tích, có độ chính xác đến ± 0,0001 g.
  \item Bếp phân hủy mẫu, có điều khiển nhiệt độ.
  \item Rây, có đường kính lỗ 0,5 mm; 1 mm.
  \item Thiết bị chưng cất Kjeldahl.
  \item Buret dung tích 50ml.
  \item Thiết bị chưng cất gồm: \begin{itemize}
    \item Bình kendan 500 - 800 ml hoặc bình cầu đáy tròn dung tích 1000 ml;
    \item Đầu nút thủy tinh hình cầu để ngăn tia bắn;
    \item Phễu nhỏ giọt dung tích 100 ml ;
    \item Ống sinh hàn làm nguội bằng nước;
    \item Bình hứng: bình tam giác hoặc cốc thủy tinh dung tích 250 ml;
    \item Bi thủy tinh đường kính 2 - 3 mm.
  \end{itemize}
\end{itemize}

\subsubsection{Quy trình phân tích}

a. Kiểm tra thiết bị chưng cất
\begin{itemize}
  \item Kiểm tra độ kín của hệ thống chưng cất bằng dung dịch tiêu chuẩn.
  \item Trước khi chưng cất mẫu phải kiểm tra thiết bị Kjeldhal bằng cách chưng cất 50 mL dung dịch tiêu chuẩn ammonium sulfate 0,5 mg N/mL với kiềm.
  Chuẩn độ lượng nitơ trong bình hứng hết 17,85 mL ± 0,1 mL dung dịch chuẩn độ sulfuric acid 0,1 N là đạt yêu cầu. 
  \begin{itemize}
    \item Nếu ít hơn là do thiết bị chưng cất bị hở.
    \item Nếu lớn hơn có thể là do bị bắn kiềm từ bình cất hoặc do thiết bị không sạch, cần khắc phục và kiểm tra lại.
  \end{itemize}
\end{itemize}

b. Chuẩn bị hoá chất
\begin{itemize}
  \item Dung dịch sodium hydroxide (\(\mathrm{NaOH}\)) 40\%: Cân chính xác 400 g \(\mathrm{NaOH}\) tinh thể, hòa tan trong khoảng 500 mL nước cất, để nguội và thêm nước đến vạch định mức 1000 mL.
  \item Phenolphthalein 1\%: Hòa tan 1 g phenolphthalein trong 100 mL cồn 96\%.
  \item Dung dịch chỉ thị màu hỗn hợp: Hoà tan 0,1g methyl đỏ tinh thể với khoảng 50ml rược etylic 95\% trong cốc dung tích 100ml, sau đó thêm 0,5g bromocresol xanh tinh thể.
  Đưa toàn bộ hỗn hợp vào bình định mức và lắc đều. Bảo quản tại 20oC, trong lọ nâu để tránh ánh sáng trực tiếp.
  \item Hỗn hợp xúc tác: Cân 1000g \(\mathrm{K_2SO_4}\) và 100g \(\mathrm{CuSO_4\cdot 2H_2O}\), trộn đều và nghiền nhỏ.
  \item Dung dịch boric acid 5\%: Cân chính xác 50 g boric acid tinh thể hoà tan trong 900 mL nước nóng. Sau khi nguội, chuyển vào bình định mức 1000mL thêm nước đến vạch định mức và lắc đều.
  \item Dung dịch ammonium sulfate 0,5 mg N/mL: Cân chính xác 2,3571g ammonium sulfate \(\mathrm{(NH_4)_2SO_4}\) (đã sấy ở 100oC trong 2h, để trong bình hút ẩm) vào cốc có dung tích 1000 mL, thêm 400ml nước, khuấy đều cho đến khi tan hết và chuyển vào bình định mức 1000 mL, thêm nước đến vạch định mức và lắc đều.
  Dung dịch được bảo quản kín tại 20oC.
\end{itemize}

c. Quy trình phân tích

c.1 Phân huỷ mẫu

\begin{itemize}
  \item Cân chính xác khoảng 0,5000g mẫu thử (±0.0002g) vào bình tam giác 100ml.
  \item Thêm 0,25g hỗn hợp xúc tác \(\mathrm{K_2SO_4:CuSO_4}\) (10:1), tẩm ướt bằng nước cất.
  \item Thêm 5ml \(\mathrm{H_2SO_4}\) đặc, đậy miệng tam giác bằng phễu nhỏ.
  \item Gia nhiệt trên bếp ở 250-300oC trong khoảng 4h-6h. Lấy ra và để nguội đến nhiệt độ phòng.
  \item Tiến hành phân hủy mẫu trắng song song với mẫu thử, với lượng thuốc thử và điều kiện tương tự mẫu thử, nhưng không có mẫu thử.
\end{itemize}

Phương pháp này áp dụng với mẫu phân bón không chứa nitơ dạng nitrat.

c.2 Chưng cất

\begin{itemize}
  \item Sau khi phân hủy mẫu và để nguội, chuyển toàn bộ dung dịch đã chuẩn bị ở (c.1) vào bình chưng cất  của thiết bị Kjeldahl.
  \item Thêm 40 mL \(\mathrm{NaOH}\) 40\% và khoảng 1 mL chỉ thị phenolphthalein 1\% vào bình chưng cất.
  \item Chuẩn bị bình hứng dung tích 250 mL có chứa 50 mL dung dịch acid boric 5\% đã có hỗn hợp chỉ thị (đuôi ống sinh hàn phải ngập trong dung dịch acid boric khoảng 2 mm).
  \item Tiến hành tách amoni bằng thiết bị chưng cất Kjeldahl. Kết thúc quá trình chưng cất khi hết amoni (khi dung dịch ngưng khoảng 150 mL).
  \item Hạ thấp bình hứng, tia rửa đuôi ống sinh hàn vào bình hứng và để nguội.
\end{itemize}

c.3 Chuẩn độ

\begin{itemize}
  \item Chuẩn độ ammonium tetraborate sau khi chưng cất bằng dung dịch chuẩn độ sulfuric acid 0,1N, lắc liên tục cho đến khi chuyển màu đột ngột từ xanh sang đỏ nhạt. 
  \item Ghi lại thể tích acid sulfuric đã tiêu tốn.
\end{itemize}

d. Tính toán kết quả
Tổng hàm lượng nitơ (N), được tính bằng phần trăm khối lượng theo công thức sau:
\[
  \%\mathrm{N} = \frac{(V_2 - V_1) \times N_{\mathrm{H_2SO_4}} \times 14{,}01}{m \times 1000} \times 100
\]

Trong đó:
\begin{itemize}
  \item $V_1$ là thể tích dung dịch sulfuric acid chuẩn dùng để chuẩn độ mẫu trắng (mL);
  \item $V_2$ là thể tích dung dịch sulfuric acid chuẩn dùng để chuẩn độ mẫu phân tích (mL);
  \item $N_{\mathrm{H_2SO_4}}$ là nồng độ của dung dịch sulfuric acid dùng để chuẩn độ (N);
  \item $m$ là khối lượng mẫu cân (g);
  \item $14{,}01$ là khối lượng mol của nitơ (g/mol).
\end{itemize}

\subsection{Xác định kali hữu hiệu}
\subsubsection{Nguyên tắc}
Nguyên tắc của phương pháp là chiết kali trong mẫu phân bón bằng dung dịch \(\mathrm{HCl}\), sau đó xác định nồng độ ion kali trong dịch chiết bằng máy quang kế ngọn lửa.
Khi dung dịch mẫu được đưa vào ngọn lửa, năng lượng nhiệt của ngọn lửa kích thích một phần nguyên tử kali lên trạng thái năng lượng cao hơn thành nguyên tử kali tự do.
\[
  \mathrm{K} + \text{năng lượng nhiệt} \longrightarrow \mathrm{K}^{*}
\]
Trong đó, $\mathrm{K}^{*}$ là nguyên tử kali ở trạng thái kích thích.

Khi các nguyên tử trở về trạng thái năng lượng ban đầu, chúng phát ra bức xạ đặc trưng của kali. 
\[
  \mathrm{K}^{*} \longrightarrow \mathrm{K} + h\nu
\]
Bức xạ $h\nu$ có bước sóng đặc trưng của kali và được thiết bị ghi nhận để định lượng.

Cường độ bức xạ phát ra được máy quang kế ngọn lửa ghi nhận và so sánh với đường chuẩn dung dịch kali để xác định nồng độ kali trong mẫu.

\subsubsection{Hóa chất, thiết bị và dụng cụ}
+ Hóa chất và thuốc thử:
\begin{itemize}
  \item Hydrochloric acid (\(\mathrm{HCl}\)) đậm đặc (d = 1,184)
  \item Cesium chloride (\(\mathrm{CsCl}\)) tinh thể.
  \item Dung dịch chuẩn kali gốc 1000 mg/L.
\end{itemize}

+ Thiết bị và dụng cụ:
\begin{itemize}
  \item Các dụng cụ thông thường trong phòng thí nghiệm.
  \item Cân phân tích, có độ chính xác đến ± 0,0001 g.
  \item Máy quang kế ngọn lửa.
  \item Rây, có đường kính lỗ 1 mm.
  \item Bếp khuấy.
\end{itemize}

\subsubsection{Quy trình phân tích}

a. Chuẩn bị hoá chất
\begin{itemize}
  \item Dung dịch hydrochloric acid 0,1N: Lấy 83,4 mL hydrochloric acid đậm đặc hòa tan với khoảng 300 mL nước trong bình định mức dung tích 1000 mL. 
  Thêm nước đến vạch định mức và lắc đều.
  \item Dung dịch hydrochloric acid 0,05N: Lấy 50 mL dung dịch hydrochloric acid 0,1N định mức 1000 mL. Ta thu được dung dịch hydrochloric acid 0,05N.
  \item Dung dịch cesium chloride 25 g/L: Hòa tan 25 g cesium chloride tinh thể trong 50 mL hydrochloric acid đậm đặc và 450 mL nước trong bình định mức dung tích 1000 mL. 
  Thêm nước đến vạch định mức và lắc đều.
  \item Dung dịch potassium chuẩn 100 mg/L: Dùng pipet hút chính xác 25 mL dung dịch chuẩn potassium gốc 1000 mg/L cho vào bình định mức dung tích 250 mL, thêm dung dịch hydrochloric acid 0,05N tới vạch định mức, lắc đều.
\end{itemize}

b. Phân huỷ mẫu
\begin{itemize}
  \item Cân chính xác khoảng 0,5000g - 1,0000g mẫu thử (±0,0001 g) đã được nghiền nhỏ, chuyển vào bình tam giác 250 mL.
  \item Thêm vào đó 100 mL dung dịch \(\mathrm{HCl}\) 0,05N và khuấy cho đến khi tan hết mẫu.
  \item Chuyển vào bình định mức 250 mL, thêm \(\mathrm{HCl}\) 0,05N đến vạch định mức, lắc đều và lọc dung dịch qua giấy lọc khô.
  \item Dung dịch đã lọc được dùng để xác định \(\mathrm{K_2O}\) hữu hiệu (dung dịch A).
  \item Tiến hành phân hủy mẫu trắng song song với mẫu thử, với lượng thuốc thử và điều kiện tương tự mẫu thử, nhưng không có mẫu thử.
\end{itemize}

c. Chuẩn bị mẫu thử
\begin{itemize}
  \item Dùng pipet hút 10 mL dung dịch A đã được chuẩn bị ở trên vào bình định mức dung tích 100 mL, và định mức tới vạch bằng \(\mathrm{HCl}\) 0,05N.
  \item Thêm dung dịch cesium chloride 25 g/L vào dung dịch A trước khi đo với tỉ lệ về thể tích dung dịch cesium chloride 25 g/L và dung dịch A là 1:1.
\end{itemize}

d. Kiểm tra máy quang kế ngọn lửa
\begin{itemize}
  \item Kiểm tra theo hướng dẫn của nhà sản xuất.
  \item Khởi động máy trước khi sử dụng 15 phút để ổn định ngọn lửa.
  \item Thiết lập các điều kiện và thông số đo theo hướng dẫn của nhà sản xuất.
\end{itemize}

\noindent\textbf{e. Xây dựng đường chuẩn}\par

\begin{table}[H]
  \centering
  \small
  \setlength{\tabcolsep}{2pt}
  \renewcommand{\arraystretch}{1.3}
  \caption{Chuẩn bị dãy dung dịch chuẩn kali để xây dựng đường chuẩn}
  \label{tab:duong-chuan-kali}
  \begin{tabular}{|>{\raggedright\arraybackslash}p{0.42\textwidth}|*{6}{>{\centering\arraybackslash}p{0.08\textwidth}|}}
    \hline
    \centering\textbf{Số hiệu bình} & \textbf{S0} & \textbf{S1} & \textbf{S2} & \textbf{S3} & \textbf{S4} & \textbf{S5} \\
    \hline
    Thể tích dung dịch chuẩn kali 100 mg/L lấy vào mỗi bình (mL) & 0 & 5 & 10 & 15 & 20 & 25 \\
    \hline
    Thể tích dung dịch \(\mathrm{CsCl}\) 25 g/L thêm vào mỗi bình (mL) & 50 & 50 & 50 & 50 & 50 & 50 \\
    \hline
    Thể tích dung dịch \(\mathrm{HCl}\) 0,05 N thêm đến vạch định mức (mL) & 50 & 47,5 & 45 & 40 & 30 & 20 \\
    \hline
    Thể tích định mức (mL) & 100 & 100 & 100 & 100 & 100 & 100 \\
    \hline
    Nồng độ dung dịch chuẩn kali (mg/L) & 0,00 & 5,0 & 10,0 & 15,0 & 20,0 & 25,0 \\
    \hline
  \end{tabular}
\end{table}


Thực hiện xây dựng đường chuẩn theo các bước sau:
\begin{itemize}
  \item Dùng pipet lấy lần lượt các thể tích dung dịch chuẩn kali 100 mg/L vào các bình định mức S0--S5 theo Bảng~\ref{tab:duong-chuan-kali}.
  \item Thêm 50 mL dung dịch \(\mathrm{CsCl}\) 25 g/L vào mỗi bình.
  \item Thêm dung dịch \(\mathrm{HCl}\) 0,05 N đến vạch định mức 100 mL theo thể tích tương ứng trong bảng.
  \item Đậy nút bình, lắc đều để thu được dãy chuẩn có nồng độ từ 0,00 đến 25,0 mg/L.
  \item Đo cường độ phát xạ của từng dung dịch chuẩn trên máy quang kế ngọn lửa và lập đường chuẩn biểu diễn mối quan hệ giữa cường độ phát xạ và nồng độ kali.
  \item Dựng đồ thị với trục hoành là nồng độ kali (mg/L) và trục tung là cường độ phát xạ.
  \[
    y = ax + b
  \]
  \item Dựng lại dãy chuẩn tối thiểu mỗi tuần một lần; sau khoảng 10 mẫu đo, kiểm tra lại đường chuẩn và hiệu chỉnh máy nếu có sai lệch.
\end{itemize}

\subsubsection{Kết quả xây dựng đường chuẩn}
Kết quả đo cường độ phát xạ của dãy chuẩn được ghi nhận theo Bảng~\ref{tab:ket-qua-duong-chuan-kali}. Dùng các giá trị này để hồi quy tuyến tính, trong đó $x$ là nồng độ kali và $y$ là cường độ phát xạ.

\begin{table}[H]
  \centering
  \small
  \renewcommand{\arraystretch}{1.3}
  \caption{Bảng ghi nhận kết quả xây dựng đường chuẩn kali}
  \label{tab:ket-qua-duong-chuan-kali}
  \resizebox{\textwidth}{!}{%
  \begin{tabular}{|>{\centering\arraybackslash}p{0.14\textwidth}|*{6}{>{\centering\arraybackslash}p{0.12\textwidth}|}}
    \hline
    \textbf{Số hiệu bình} & \textbf{S0} & \textbf{S1} & \textbf{S2} & \textbf{S3} & \textbf{S4} & \textbf{S5} \\
    \hline
    Nồng độ K (mg/L) & 0,00 & 5,0 & 10,0 & 15,0 & 20,0 & 25,0 \\
    \hline
    Cường độ phát xạ & 0 & 5 & 10 & 16 & 21 & 25 \\
    \hline
  \end{tabular}
  }
\end{table}


\begin{figure}[H]
  \centering
  \begin{tikzpicture}
    \begin{axis}[
      width=0.82\textwidth,
      height=0.52\textwidth,
      xlabel={Nồng độ kali (mg/L)},
      ylabel={Cường độ phát xạ},
      xmin=0, xmax=26,
      ymin=0, ymax=27,
      xtick={0,5,10,15,20,25},
      ytick={0,5,10,15,20,25},
      grid=both,
      major grid style={gray!35},
      minor grid style={gray!15},
      legend style={at={(0.97,0.04)},anchor=south east,font=\small},
      legend cell align={left}
    ]
      \addplot[
        only marks,
        mark=*,
        mark size=2.3pt,
        blue!75!black
      ] coordinates {(0,0) (5,5) (10,10) (15,16) (20,21) (25,25)};
      \addlegendentry{Số liệu chuẩn}
      \addplot[red!75!black, very thick, domain=0:25, samples=100] {1.022857142857*x + 0.047619047619};
      \addlegendentry{Hồi quy tuyến tính}
      \node[
        anchor=north west,
        draw=black,
        fill=white,
        fill opacity=0.92,
        text opacity=1,
        align=left,
        font=\small
      ] at (axis cs:1,25.8) {$y = 1{,}0229x + 0{,}0476$\\$R^2 = 0{,}9976$};
    \end{axis}
  \end{tikzpicture}
  \caption{Đường hồi quy tuyến tính của dãy chuẩn kali}
  \label{fig:duong-hoi-quy-kali}
\end{figure}


\noindent\textbf{Nhận xét đường chuẩn:} Cường độ phát xạ phải tăng theo nồng độ kali trong khoảng 0,00--25,0 mg/L. Đường chuẩn được chấp nhận khi mối quan hệ tuyến tính phù hợp với yêu cầu kiểm soát nội bộ của phòng thí nghiệm; nếu điểm chuẩn có độ lệch bất thường hoặc hệ số tương quan không đạt yêu cầu, cần kiểm tra lại dung dịch chuẩn, thao tác pha loãng và điều kiện vận hành máy.

\noindent\textbf{f. Tính toán kết quả}\par
Nồng độ kali trong dung dịch mẫu được nội suy từ phương trình đường chuẩn $I=aC+b$:
\[
  C_{\mathrm{K}} = \frac{I-b}{a}
\]
Trong đó $I$ là cường độ phát xạ của dung dịch mẫu, $C_{\mathrm{K}}$ là nồng độ kali nội suy (mg/L), $a=1{,}0229$ là hệ số góc và $b=0{,}0476$ là hệ số chặn của đường hồi quy.

Hàm lượng kali hữu hiệu, tính theo phần trăm khối lượng, được xác định theo công thức:
\[
  \%\mathrm{K} = \frac{\left(C_{\mathrm{K}}-C_{\mathrm{blank}}\right) \times f \times V}{m \times 10\,000}
\]

Hàm lượng kali quy đổi về \(\mathrm{K_2O}\) được tính như sau:
\[
  \%\mathrm{K_2O} = \%\mathrm{K} \times 1{,}205
\]

Trong đó:
\begin{itemize}
  \item $C_{\mathrm{K}}$ là nồng độ kali trong dung dịch mẫu, nội suy từ đường chuẩn (mg/L);
  \item $C_{\mathrm{blank}}$ là nồng độ kali trong dung dịch mẫu trắng (mg/L);
  \item $f$ là hệ số pha loãng;
  \item $V$ là thể tích dung dịch A sau khi chiết (mL);
  \item $m$ là khối lượng mẫu cân (g);
  \item $1{,}205$ là hệ số quy đổi từ K sang \(\mathrm{K_2O}\).
\end{itemize}

\subsection{Xác định lân hữu hiệu}
\subsubsection{Nguyên tắc}
Đối với kiểm định phân lân, chỉ tiêu thường được quan tâm là hàm lượng phốt pho hữu hiệu, được biểu thị dưới dạng \(\mathrm{P_2O_5}\).
Phospho hữu hiệu trong mẫu phân bón được chiết bằng dung dịch Petermann, là dung dịch ammonium citrate trong môi trường kiềm.
Ion citrate trong dung dịch Petermann có khả năng tạo phức với ion kim loại, nhờ đó thúc đẩy quá trình hoà tan các dạng muối phosphat có trong phân bón và chuyển phosphat cần xác định vào trong dung dịch.
Quá trình hòa tan phosphat có thể được biểu diễn khái quát như sau:
\[
  \mathrm{M_3(PO_4)}_n \rightleftharpoons 3\mathrm{M}^{n+} + n\mathrm{PO_4^{3-}}
\]
Các ion kim loại giải phóng tiếp tục tạo phức với citrate, làm dịch chuyển cân bằng theo hướng hòa tan:
\[
  \mathrm{M}^{n+} + \mathrm{Cit}^{3-} \rightleftharpoons [\mathrm{M(Cit)}]^{(n-3)+}
\]
Trong đó, $\mathrm{M^{n+}}$ là ion kim loại và $\mathrm{Cit^{3-}}$ là ion citrate.

Sau khi thu được dịch chiết trong, lấy một thể tích xác định của dịch này để định lượng phospho. 
Ion phosphate trong môi trường axit được cho phản ứng với ammonium molybdate, tạo kết tủa màu vàng ammonium phosphomolybdate:
\[
  \mathrm{PO_4^{3-} + 3NH_4^+ + 12MoO_4^{2-} + 24H^+
  \longrightarrow (NH_4)_3PO_4 \cdot 12MoO_3 \downarrow + 12H_2O}
\]
Kết tủa ammonium phosphomolybdate được lọc và rửa sạch nhằm loại bỏ lượng axit và thuốc thử còn dư. Sau đó, kết tủa được hòa tan bằng dung dịch \(\mathrm{NaOH}\) chuẩn.
\[
  \mathrm{(NH_4)_3PO_4 \cdot 12MoO_3 + 23NaOH
  \longrightarrow 11Na_2MoO_4 + (NH_4)_2MoO_4 + NaNH_4HPO_4 + 11H_2O}
\]
Phần \(\mathrm{NaOH}\) dư được chuẩn độ bằng dung dịch \(\mathrm{H_2SO_4}\) chuẩn theo phản ứng trung hòa:
\[
  \mathrm{2NaOH + H_2SO_4 \longrightarrow Na_2SO_4 + 2H_2O}
\]

Sau khi kết tủa đã hòa tan hoàn toàn, lượng \(\mathrm{NaOH}\) còn dư được chuẩn độ bằng dung dịch axit chuẩn \(\mathrm{H_2SO_4}\).
Từ đó tính được lượng \(\mathrm{NaOH}\) đã phản ứng với kết tủa và tính được hàm lượng phospho, sau đó quy đổi kết quả về \(\%\mathrm{P_2O_5}\) trong mẫu phân bón.

\subsubsection{Hóa chất, thiết bị và dụng cụ}
+ Hoá chất và thuốc thử
\begin{itemize}
  \item Dung dịch sodium hydroxide (\(\mathrm{NaOH}\)) 0,1 N;
  \item Dung dịch citric acid đậm đặc (d=1,42);
  \item Dung dịch sulfuric acid (\(\mathrm{H_2SO_4}\)) 0,1 N;
  \item Phenolphthalein (\(\mathrm{C_{20}H_{12}O_4}\)) dạng tinh thể;
  \item Metyl đỏ (\(\mathrm{C_{15}H_{15}N_3O_2}\)) dạng tinh thể;
  \item Xanh metylen (\(\mathrm{C_{16}H_{18}ClN_3S}\)) dạng tinh thể;
  \item Potassium nitrate (\(\mathrm{KNO_3}\)) dạng tinh thể;
  \item Ammonium molybdate (\(\mathrm{(NH_4)_6Mo_7O_{24}\cdot 4H_2O}\)) dạng tinh thể;
  \item Dung dịch ammonium hydroxide (\(\mathrm{NH_4OH}\)) 10\%;
  \item Dung dịch diammonium phosphate 10\%;
\end{itemize}


+ Thiết bị và dụng cụ 
\begin{itemize}
  \item Các dụng cụ thông thường trong phòng thí nghiệm.
  \item Cân phân tích, có độ chính xác đến ± 0,0001 g.
  \item Bếp từ.
  \item Rây có đường kính lỗ 0,5 mm và 1 mm.
  \item Burrette 50 mL.
\end{itemize}

\subsubsection{Cách tiến hành}
\noindent\textbf{a. Chuẩn bị thuốc thử}
\begin{itemize}
  \item Phenolphthalein 1\%: Hòa tan 1 g phenolphthalein trong 100 mL cồn 96\%.
  \item Dung dịch chỉ thị hỗn hợp: Hòa tan 0,10 g metyl đỏ trong khoảng 50 mL ethanol 95\%, thêm 0,05 g xanh metylen và lắc đến tan hoàn toàn. 
  Định mức bằng ethanol 95\% đến 100 mL, lắc đều và bảo quản trong lọ nâu ở $20\,^{\circ}\mathrm{C}$.
  \item Dung dịch potassium nitrate (\(\mathrm{KNO_3}\)) 1\%: Hòa tan 1 g \(\mathrm{KNO_3}\) trong 100 mL nước cất.
  \item Thuốc thử ammonium molybdate 3\%: Cho 832 mL nước vào cốc lần lượt thêm 27 g ammonium nitrate, 26,3 g citric acid và 34 g ammonium molybdate; khuấy đến tan. 
  Thêm 28,7 mL \(\mathrm{HNO_3}\) đậm đặc ($d=1{,}4$), đun sôi, sau đó thêm 2 giọt dung dịch diammonium phosphate 10\%.
  Khuấy đều, để nguội, để yên qua đêm, lọc qua giấy lọc và bảo quản trong chai thủy tinh màu nâu.
  \item Chuẩn hóa dung dịch \(\mathrm{NH_4OH}\) 10\% để pha Petermann: Lấy 10 mL dung dịch \(\mathrm{NH_4OH}\) 10\% vào bình định mức 500 mL đã chứa khoảng 400--450 mL nước; thêm nước đến vạch và lắc đều.
  Lấy 25 mL dung dịch đã pha loãng vào bình tam giác 250 mL chứa 25--50 mL nước, thêm 2--3 giọt metyl đỏ và chuẩn độ bằng \(\mathrm{H_2SO_4}\) 0,1 N đến màu hồng; ghi thể tích axit tiêu tốn là $V$ (mL). Thể tích \(\mathrm{NH_4OH}\) 10\% cần lấy để pha 1 L Petermann được tính theo công thức:
  \[
    X = \frac{42 \times 10 \times 25}{V \times 0{,}0014 \times 500}
  \]
  trong đó $X$ là thể tích \(\mathrm{NH_4OH}\) 10\% cần lấy (mL) và $0{,}0014$ là lượng nitơ tương ứng với 1 mL \(\mathrm{H_2SO_4}\) 0,1 N (g).
  \item Dung dịch Petermann: Chuyển thể tích $X$ dung dịch \(\mathrm{NH_4OH}\) 10\% vào bình định mức 1 L. Hòa tan 173 g citric acid trong 200--250 mL nước nóng, để nguội về nhiệt độ phòng rồi rót từ từ vào bình chứa \(\mathrm{NH_4OH}\), đồng thời làm mát bình. Tráng cốc và phễu bằng nước, chuyển nước tráng vào bình, để nguội hoàn toàn, thêm nước đến vạch 1 L và lắc kỹ. Để dung dịch ổn định trong 2 ngày trước khi sử dụng.
\end{itemize}

\noindent\textbf{b. Quy trình chiết phospho hữu hiệu}
\begin{itemize}
  \item Cân khoảng 2,0--2,5 g mẫu đã nghiền nhỏ vào bình tam giác 250 mL.
  \item Thêm 25 mL dung dịch Petermann, nghiền hoặc khuấy mẫu trong vài phút rồi để cặn lắng.
  \item Gạn phần dịch chiết vào bình định mức 250 mL qua phễu thủy tinh.
  \item Lặp lại thao tác chiết thêm ba lần, mỗi lần dùng 25 mL dung dịch Petermann.
  \item Chuyển toàn bộ cặn còn lại trong cối và trên phễu vào bình định mức; dùng tổng cộng 200 mL dung dịch Petermann cho quá trình chiết.
  \item Đặt bình trong bếp cách thủy ở $60\,^{\circ}\mathrm{C} \pm 1\,^{\circ}\mathrm{C}$ trong 15 phút.
  \item Để nguội, thêm dung dịch Petermann đến vạch 250 mL, lắc đều và lọc để thu dịch chiết trong.
  \item Chuẩn bị mẫu trắng song song, dùng cùng lượng thuốc thử và điều kiện thao tác nhưng không có mẫu phân bón.
\end{itemize}

\noindent\textbf{c. Quy trình hòa tan mẫu}
\begin{itemize}
  \item Pipet chính xác 50 mL dịch chiết trong vào cốc dung tích 250 mL.
  \item Thêm 5 mL \(\mathrm{HNO_3}\) đậm đặc, đun sôi nhẹ trong 5 phút rồi để nguội.
  \item Nhỏ từ từ dung dịch \(\mathrm{NH_4OH}\) đến khi bắt đầu xuất hiện kết tủa.
  \item Thêm \(\mathrm{HNO_3}\) 6 N để hòa tan hoàn toàn kết tủa vừa tạo thành.
\end{itemize}

\noindent\textbf{d. Tiến hành thử nghiệm -- tạo kết tủa}
\begin{itemize}
  \item Đun sôi dung dịch đã xử lý trên bếp điện, đồng thời khuấy nhẹ.
  \item Thêm từ từ 50 mL thuốc thử amoni molipdat 3\% vào dung dịch.
  \item Tiếp tục đun sôi trong 3 phút, sau đó để kết tủa lắng trong khoảng 30 phút đến 2 giờ.
  \item Lọc kết tủa qua giấy lọc băng vàng.
  \item Rửa kết tủa bằng dung dịch \(\mathrm{KNO_3}\) 1\% đến khi loại hết axit; kiểm tra nước rửa bằng phenolphthalein theo SOP áp dụng.
\end{itemize}

\noindent\textbf{e. Quy trình chuẩn độ}
\begin{itemize}
  \item Chuyển giấy lọc có kết tủa trở lại cốc đã dùng để tạo kết tủa.
  \item Dùng buret thêm từ từ dung dịch \(\mathrm{NaOH}\) 0,1 N để hòa tan hoàn toàn kết tủa.
  \item Thêm dư 5 mL \(\mathrm{NaOH}\) 0,1 N sau khi kết tủa đã tan hoàn toàn và ghi tổng thể tích \(\mathrm{NaOH}\) đã dùng là $V_1$ (mL).
  \item Rửa thành cốc bằng bình tia và thêm 3--5 giọt phenolphthalein.
  \item Chuẩn độ lượng \(\mathrm{NaOH}\) dư bằng dung dịch \(\mathrm{H_2SO_4}\) 0,1 N đến khi màu hồng vừa mất; ghi thể tích \(\mathrm{H_2SO_4}\) đã dùng là $V_2$ (mL).
\end{itemize}

\subsubsection{Tính toán kết quả}
Hàm lượng \(\mathrm{P_2O_5}\) hữu hiệu, tính theo phần trăm khối lượng, được xác định theo công thức:
\[
  \%\mathrm{P_2O_5} = \frac{(V_1-V_2) \times 0{,}000309 \times 250}{m \times 50} \times 100
\]

Trong đó:
\begin{itemize}
  \item $V_1$ là thể tích \(\mathrm{NaOH}\) 0,1 N dùng để hòa tan kết tủa của mẫu thử (mL);
  \item $V_2$ là thể tích \(\mathrm{H_2SO_4}\) 0,1 N dùng để chuẩn độ lượng kiềm dư (mL);
  \item $0{,}000309$ là lượng \(\mathrm{P_2O_5}\) tương ứng với 1 mL \(\mathrm{NaOH}\) 0,1 N (g);
  \item $m$ là khối lượng mẫu cân (g).
\end{itemize}

\subsection{Xác định carbon tổng}
\subsubsection{Nguyên tắc}

Phương pháp Walkley--Black được dùng để xác định hàm lượng carbon hữu cơ trong đất hoặc vật liệu hữu cơ bằng phương pháp oxy hóa ướt. 
Nguyên tắc chính là dùng potassium dichromate (\(\mathrm{K_2Cr_2O_7}\)) dư trong môi trường \(\mathrm{H_2SO_4}\) đậm đặc để oxy hóa carbon hữu cơ trong mẫu.

\noindent\textbf{a. Oxy hóa carbon hữu cơ}

Khi cho \(\mathrm{K_2Cr_2O_7}\) và \(\mathrm{H_2SO_4}\) đậm đặc vào mẫu, carbon hữu cơ bị oxy hóa thành \(\mathrm{CO_2}\), còn crom(VI) trong ion \(\mathrm{Cr_2O_7^{2-}}\) bị khử thành \(\mathrm{Cr^{3+}}\).

Có thể biểu diễn tổng quát:
\[
  3\mathrm{C} + 2\mathrm{Cr_2O_7^{2-}} + 16\mathrm{H^+}
  \longrightarrow 3\mathrm{CO_2} + 4\mathrm{Cr^{3+}} + 8\mathrm{H_2O}
\]

Hay xét riêng bán phản ứng của dicromat:
\[
  \mathrm{Cr_2O_7^{2-}} + 14\mathrm{H^+} + 6e^-
  \longrightarrow 2\mathrm{Cr^{3+}} + 7\mathrm{H_2O}
\]

Trong đó, lượng \(\mathrm{Cr_2O_7^{2-}}\) bị tiêu hao tỷ lệ với lượng carbon hữu cơ bị oxy hóa.

\noindent\textbf{b. Chuẩn độ lượng dicromat còn dư}

Sau khi phản ứng oxy hóa hoàn tất, lượng \(\mathrm{K_2Cr_2O_7}\) còn dư được xác định bằng chuẩn độ ngược với dung dịch \(\mathrm{Fe^{2+}}\).

Phản ứng chuẩn độ:
\[
  \mathrm{Cr_2O_7^{2-}} + 6\mathrm{Fe^{2+}} + 14\mathrm{H^+}
  \longrightarrow 2\mathrm{Cr^{3+}} + 6\mathrm{Fe^{3+}} + 7\mathrm{H_2O}
\]

Tiến hành đồng thời một mẫu trắng không chứa mẫu phân tích.
\begin{itemize}
  \item Mẫu trắng cho biết tổng lượng \(\mathrm{K_2Cr_2O_7}\) ban đầu.
  \item Mẫu kiểm định cho biết lượng \(\mathrm{K_2Cr_2O_7}\) còn dư sau khi oxy hóa carbon.
\end{itemize}

Do đó:
\[
  \boxed{\text{Dicromat đã phản ứng} = \text{Dicromat mẫu trắng} - \text{Dicromat mẫu kiểm định}}
\]

Từ lượng dicromat đã tiêu hao có thể tính được lượng carbon hữu cơ trong mẫu.

\noindent\textbf{c. Nhận biết điểm kết thúc}

Chất chỉ thị được sử dụng là ferroin, một phức chất của \(\mathrm{Fe^{2+}}\) với phenanthrolin.

Khi chuẩn độ bằng dung dịch \(\mathrm{Fe^{2+}}\), đến điểm tương đương toàn bộ \(\mathrm{Cr_2O_7^{2-}}\) dư đã bị khử.
Một lượng rất nhỏ \(\mathrm{Fe^{2+}}\) dư làm chỉ thị chuyển màu, từ đó xác định điểm kết thúc chuẩn độ.

\noindent\textbf{d. Điểm đặc biệt của Walkley--Black}

Phương pháp Walkley--Black không oxy hóa hoàn toàn 100\% carbon hữu cơ trong mọi loại mẫu. 
Phương pháp cổ điển thường giả định khoảng 75\% carbon hữu cơ được oxy hóa, vì vậy thường sử dụng hệ số hiệu chỉnh:
\[
  \frac{100}{75} \approx 1{,}33
\]

\subsubsection{Hóa chất, thiết bị và dụng cụ}
\noindent\textbf{Hóa chất và thuốc thử}
\begin{itemize}
  \item Sulfuric acid đậm đặc, \(\mathrm{H_2SO_4}\) ($d = 1{,}84$).
  \item Phosphoric acid, \(\mathrm{H_3PO_4}\) 85\%.
  \item Potassium dichromate, \(\mathrm{K_2Cr_2O_7}\).
  \item Mohr's salt, \(\mathrm{(NH_4)_2Fe(SO_4)_2\cdot 6H_2O}\).
  \item Iron(II) sulfate heptahydrate, \(\mathrm{FeSO_4\cdot 7H_2O}\).
  \item o-Phenanthroline monohydrate, \(\mathrm{C_{12}H_8N_2\cdot H_2O}\).
\end{itemize}

\noindent\textbf{Thiết bị và dụng cụ}
\begin{itemize}
  \item Dụng cụ thông thường của phòng thí nghiệm.
  \item Cân phân tích có độ chính xác đến $\pm0{,}0001$ g.
  \item Tủ sấy có nhiệt độ làm việc đến $200\,^{\circ}\mathrm{C}$ và độ chính xác $\pm2\,^{\circ}\mathrm{C}$.
  \item Rây có đường kính lỗ 0,2 mm.
  \item Buret, dung tích 50 mL và độ chính xác 0,1 mL.
\end{itemize}

\subsubsection{Cách tiến hành}

\noindent\textbf{a. Chuẩn bị thuốc thử}
\begin{itemize}
  \item \textbf{Dung dịch tiêu chuẩn \(\mathrm{K_2Cr_2O_7}\) M/6:} Cân 49,040 g \(\mathrm{K_2Cr_2O_7}\) đã sấy khô ở $105\,^{\circ}\mathrm{C}$ trong 2 giờ và để nguội trong bình hút ẩm. Cho hóa chất vào cốc dung tích 1000 mL, thêm khoảng 400 mL nước cất, khuấy đến tan hoàn toàn. Chuyển định lượng vào bình định mức 1000 mL, thêm nước cất đến vạch, lắc đều và bảo quản kín ở $20\,^{\circ}\mathrm{C}$.
  \item \textbf{Dung dịch Mohr's salt khoảng 0,5 M:} Cân 196 g \(\mathrm{(NH_4)_2Fe(SO_4)_2\cdot 6H_2O}\) vào cốc dung tích 1000 mL. Thêm 50 mL \(\mathrm{H_2SO_4}\) đậm đặc và 450 mL nước cất, khuấy đến khi muối tan hoàn toàn. Chuyển dung dịch vào bình định mức 1000 mL, thêm nước cất đến vạch và lắc đều. Để dung dịch lắng trong; nếu dung dịch bị đục thì lọc. Bảo quản dung dịch kín trong chai thủy tinh màu nâu ở $20\,^{\circ}\mathrm{C}$ và hạn chế sự xâm nhập của không khí.
  \item \textbf{Dung dịch chỉ thị màu ferroin (o-phenanthroline):} Cân 0,695 g iron(II) sulfate heptahydrate, \(\mathrm{FeSO_4\cdot 7H_2O}\), và 1,485 g o-phenanthroline monohydrate, \(\mathrm{C_{12}H_8N_2\cdot H_2O}\). Hòa tan đồng thời hai chất trong 100 mL nước cất, lắc đều đến khi dung dịch trong suốt; bảo quản trong chai màu nâu.
\end{itemize}

\noindent\textbf{b. Chuẩn bị mẫu}
\begin{itemize}
  \item Trộn đều mẫu phân bón.
  \item Nghiền mịn mẫu.
  \item Rây mẫu qua rây có đường kính lỗ 0,2 mm.
  \item Trộn đều phần mẫu đã rây để sử dụng cho phân tích.
\end{itemize}

\noindent\textbf{c. Oxy hóa carbon hữu cơ}
\begin{itemize}
  \item Cân chính xác khoảng 0,1000--0,2000 g mẫu vào bình tam giác chịu nhiệt dung tích 250 mL.
  \item Thêm chính xác 20,0 mL dung dịch tiêu chuẩn \(\mathrm{K_2Cr_2O_7}\) M/6 vào bình.
  \item Thêm nhanh 40 mL \(\mathrm{H_2SO_4}\) đậm đặc.
  \item Lắc nhẹ bình để trộn đều dung dịch.
  \item Đặt bình trên tấm cách nhiệt và để yên trong 60 phút.
  \item Kiểm tra màu dung dịch sau oxy hóa; nếu dung dịch có màu xanh, thực hiện lại phép thử với lượng mẫu nhỏ hơn hoặc tăng lượng dung dịch \(\mathrm{K_2Cr_2O_7}\).
\end{itemize}

\noindent\textbf{d. Chuẩn bị mẫu trắng và xử lý sau oxy hóa}
\begin{itemize}
  \item Chuẩn bị đồng thời hai mẫu trắng theo cùng quy trình nhưng không thêm mẫu phân bón.
  \item Thêm 100 mL nước cất vào từng bình sau thời gian oxy hóa.
  \item Thêm 10 mL \(\mathrm{H_3PO_4}\) 85\% vào từng bình.
  \item Để dung dịch nguội đến nhiệt độ phòng trước khi chuẩn độ.
\end{itemize}

\noindent\textbf{e. Chuẩn độ}
\begin{itemize}
  \item Thêm 0,5 mL dung dịch chỉ thị ferroin vào mẫu trắng và mẫu thử.
  \item Chuẩn độ lượng \(\mathrm{K_2Cr_2O_7}\) dư bằng dung dịch muối Mohr khoảng 0,5 M.
  \item Nhỏ từ từ từng giọt dung dịch muối Mohr khi gần đến điểm kết thúc và lắc đều bình.
  \item Dừng chuẩn độ khi dung dịch chuyển màu đột ngột từ xanh sang đỏ thẫm.
  \item Ghi thể tích dung dịch muối Mohr đã sử dụng cho mẫu trắng và mẫu thử.
\end{itemize}

\subsubsection{Tính toán kết quả}
Hàm lượng carbon hữu cơ, tính theo phần trăm khối lượng, được xác định theo công thức:
\[
  \%\mathrm{OC} = \frac{(a-b) \times N_{\mathrm{Mohr}} \times 0{,}003 \times 100}{m} \times \frac{100}{75}
\]

Hàm lượng chất hữu cơ được quy đổi theo công thức:
\[
  \%\mathrm{OM} = \%\mathrm{OC} \times 2{,}2
\]

Trong đó:
\begin{itemize}
  \item $a$ là thể tích dung dịch muối Mohr dùng để chuẩn độ mẫu trắng (mL);
  \item $b$ là thể tích dung dịch muối Mohr dùng để chuẩn độ mẫu thử (mL);
  \item $N_{\mathrm{Mohr}}$ là nồng độ đương lượng của dung dịch muối Mohr (N);
  \item $m$ là khối lượng mẫu cân (g);
  \item $0{,}003$ là khối lượng carbon tương ứng với 1 mL dung dịch chuẩn 1 N (g);
  \item $\frac{100}{75}$ là hệ số hiệu chỉnh do phương pháp oxy hóa khoảng 75\% carbon hữu cơ;
  \item $2{,}2$ là hệ số quy đổi từ carbon hữu cơ sang chất hữu cơ.
\end{itemize}

\subsection{Xác định độ ẩm}
\subsubsection{Nguyên tắc}
Độ ẩm của mẫu được xác định bằng phương pháp khối lượng, dựa trên sự giảm khối lượng của mẫu sau khi sấy ở nhiệt độ quy định đến khi đạt khối lượng không đổi. 
Phần khối lượng mất đi trong quá trình sấy được xem chủ yếu là lượng nước có trong mẫu.

Quá trình được thực hiện bằng cách cân một lượng mẫu xác định, sau đó sấy trong tủ sấy ở nhiệt độ thích hợp. 
Sau khi sấy, mẫu được làm nguội trong bình hút ẩm rồi cân lại. 
Việc sấy và cân được lặp lại cho đến khi chênh lệch giữa hai lần cân liên tiếp không đáng kể.

\subsubsection{Hóa chất, thiết bị và dụng cụ}
\noindent\textbf{Hóa chất và thuốc thử}
\begin{itemize}
  \item Không sử dụng thuốc thử trong phép xác định độ ẩm.
  \item Nước cất được sử dụng khi cần vệ sinh dụng cụ.
\end{itemize}

\noindent\textbf{Thiết bị và dụng cụ}
\begin{itemize}
  \item Dụng cụ thông thường của phòng thí nghiệm.
  \item Cân phân tích có độ chính xác đến $\pm0{,}0001$ g.
  \item Tủ sấy có thông gió, điều chỉnh được từ $50$ đến $200\,^{\circ}\mathrm{C}$ với độ chính xác $\pm2\,^{\circ}\mathrm{C}$.
  \item Chén cân bằng thủy tinh hoặc nhôm, có nắp đậy và dung tích khoảng 30 mL.
  \item Bình hút ẩm.
\end{itemize}

\subsubsection{Cách tiến hành}
\noindent\textbf{a. Chọn nhiệt độ sấy}
\begin{itemize}
  \item Sấy ở $50$--$60\,^{\circ}\mathrm{C}$ đối với phân bón kém bền nhiệt có chứa nitơ.
  \item Sấy ở $70\,^{\circ}\mathrm{C}$ đối với phân hữu cơ khoáng, phân hữu cơ sinh học hoặc phân hữu cơ vi sinh.
  \item Sấy ở $105\,^{\circ}\mathrm{C}$ đối với phân bón bền nhiệt như tecmophotphat hoặc supephotphat.
\end{itemize}

\noindent\textbf{b. Chuẩn bị chén cân}
\begin{itemize}
  \item Sấy chén cân ở $105\,^{\circ}\mathrm{C}$ trong 1 giờ.
  \item Chuyển chén cân vào bình hút ẩm và để nguội đến nhiệt độ phòng.
  \item Cân và ghi khối lượng chén cân là $m_c$.
\end{itemize}

\noindent\textbf{c. Cân mẫu trước khi sấy}
\begin{itemize}
  \item Cân khoảng 5--10 g mẫu vào chén cân đã biết khối lượng.
  \item Đậy nắp chén cân.
  \item Cân và ghi khối lượng chén có mẫu trước khi sấy là $(m_c+m_m)_{\mathrm{ts}}$.
\end{itemize}

\noindent\textbf{d. Sấy mẫu và cân đến khối lượng không đổi}
\begin{itemize}
  \item Mở nắp chén cân.
  \item Đặt chén cân chứa mẫu vào tủ sấy ở nhiệt độ đã chọn.
  \item Sấy mẫu trong 2--4 giờ.
  \item Đậy nắp chén cân sau khi kết thúc thời gian sấy.
  \item Chuyển chén cân vào bình hút ẩm và để nguội đến nhiệt độ phòng.
  \item Cân và ghi khối lượng chén có mẫu sau sấy là $(m_c+m_m)_{\mathrm{ss}}$.
  \item Lặp lại chu trình sấy, làm nguội và cân sau mỗi 30--60 phút.
  \item Kết thúc khi khối lượng mẫu không đổi.
\end{itemize}

\subsubsection{Tính toán kết quả}
Độ ẩm của mẫu phân bón, tính theo phần trăm khối lượng, được xác định theo công thức:
\[
  A\% = \frac{(m_c+m_m)_{\mathrm{ts}}-(m_c+m_m)_{\mathrm{ss}}}{(m_c+m_m)_{\mathrm{ts}}-m_c} \times 100
\]

Trong đó:
\begin{itemize}
  \item $m_c$ là khối lượng chén sau khi đã sấy ở nhiệt độ $105\,^{\circ}\mathrm{C}$, tính bằng gam;
  \item $(m_c+m_m)_{\mathrm{ts}}$ là khối lượng của mẫu và chén trước khi sấy, tính bằng gam;
  \item $(m_c+m_m)_{\mathrm{ss}}$ là khối lượng của mẫu và chén sau khi sấy, tính bằng gam.
\end{itemize}
